% Source: https://artofproblemsolving.com/community/c3h1486851_some_help
% Stable identity: aops-topic-1486851-q01
The outer edge is a circle of circumference $660$ cm, so its radius is
\[
R=\dfrac{660}{2\pi}\approx\dfrac{660\cdot7}{2\cdot22}=105\ \text{cm}.
\]
The well's radius is $75$ cm, so the width of the parapet is $105-75=30$ cm, answer $\boxed{\textbf{(A)}}$.

\begin{asy}
usepackage("amsmath");
size(5cm);
defaultpen(fontsize(10pt));
real r=75, R=105;
filldraw(circle((0,0),R),lightgray,linewidth(0.9));
filldraw(circle((0,0),r),white,linewidth(0.9));
draw((-r,0)--(r,0),blue,Arrows(6));
label("$150$",(0,0),N,blue);
draw((r,0)--(R,0),red+linewidth(0.9));
label("$w$",((r+R)/2,0),N,red);
label("outer edge $660$",R*dir(-60),SE);
label("\mbox{}",(min(currentpicture,user=true)+max(currentpicture,user=true))/2,black); // keeps the text after the figure black when LaTeX typesets the labels inline
\end{asy}

(With the exact value of $\pi$ the width is $\dfrac{330}{\pi}-75\approx30.04$ cm, still $30$ cm to the nearest centimetre.)
